\section{Rust and Python}
\label{sec:RustPythonImplementation}

IPTO includes both a Rust implementation area and Python-facing entry points, and these two parts are intentionally connected. The Rust implementation exposes IPTO into a Python environment through \texttt{pyo3} and \texttt{maturin}. 

\subsection{Rust}

The Rust code under \texttt{implementations/rust} is another implementation surface for the same conceptual model. It contains:
\begin{description}
\item[\texttt{src/repo.rs}] Backend-agnostic repository orchestration.
\vspace{.2cm}
\item[\texttt{src/backend.rs}] The backend trait that defines the implementation contract.
\vspace{.2cm}
\item[\texttt{src/backends/postgres.rs}] The PostgreSQL backend.
\vspace{.2cm}
\item[\texttt{src/backends/neo4j.rs}] The Neo4j backend.
\vspace{.2cm}
\item[\texttt{src/graphql.rs} and \texttt{src/graphql\_sdl.rs}] A Rust-side GraphQL runtime and SDL handling layer.
\vspace{.2cm}
\item[\texttt{src/lib.rs}] The Python module entry point when the \texttt{python} feature is enabled.
\end{description}

The Rust implementation is \emph{maturin-ready} and exposes a stable Python API via \texttt{pyo3}. This means the Rust code is not just an internal library. It is also the provider of the Python-visible IPTO surface.

\subsection{PyO3 bridge}

When the Rust crate is built with the \texttt{python} feature, \texttt{src/lib.rs} exposes a Python class named \texttt{PyIpto}. This class wraps the Rust repository service and backend implementations and exports repository and GraphQL operations into Python.

The Python-visible facade includes methods for:
\begin{itemize}
\item health checks and statistics,
\item unit storage and retrieval,
\item searches,
\item status changes,
\item lock handling,
\item relation and association handling,
\item attribute and tenant metadata lookup,
\item GraphQL execution and SDL inspection/configuration.
\end{itemize}

This is the meaningful system-level connection between Rust and Python: Python is used as a host environment for the Rust implementation.

\subsection{Build and packaging model}

The Rust implementation is packaged for Python through \texttt{maturin}. The Rust project includes \texttt{Cargo.toml}, \texttt{pyproject.toml} and test/benchmark code that exercise the Python bridge. Python integration tests and soak benchmarks run through the \texttt{ipto\_rust.PyIpto} module, so the PyO3 path is treated as a first-class supported interface.

\subsection{Backend abstraction}

As in the Erlang implementation, backend plurality is explicit in the Rust code. The trait \texttt{Backend} in \texttt{src/backend.rs} defines the storage contract for:
\begin{itemize}
\item unit retrieval, existence checks and storage,
\item search,
\item relations and associations,
\item locks and status changes,
\item attribute and tenant metadata,
\item template persistence hooks,
\item health and cache-control hooks.
\end{itemize}

This trait is intentionally repository-shaped rather than database-shaped. It does not expose raw SQL or driver-level concerns. That keeps \texttt{RepoService} as the main place where IPTO semantics are enforced while allowing PostgreSQL and Neo4j backends to diverge internally.

\subsection{Repository service layer}

The Rust orchestration layer lives in \texttt{src/repo.rs} as \texttt{RepoService}. It plays a role very similar to the Java \texttt{Repository} and the Erlang \texttt{ipto\_repo} module.

In particular, \texttt{RepoService} is responsible for:
\begin{itemize}
\item delegating unit reads and writes to the selected backend,
\item enforcing lifecycle rules around readonly state, locks and status transitions,
\item composing native or service-level set-operation searches,
\item exposing relation, association, attribute and tenant operations through a uniform surface,
\item parsing and applying SDL-driven configuration.
\end{itemize}

One important design choice is that Rust does not rely only on backend-native search composition. If the backend reports that a native set-expression execution path is unsupported, \texttt{RepoService} can evaluate boolean set operations in the service layer and then page the merged results. That keeps the external behavior more stable across backends even when their native capabilities differ.

Rust also supports textual and symbolic search input, but that surface is currently simpler than the Java textual-query language. The Rust implementation is converging on Java semantics, yet the Java parser remains the more complete reference point for complex textual queries.

\subsection{GraphQL runtime scope}

The Rust implementation includes its own GraphQL runtime in \texttt{src/graphql.rs}. This runtime is not yet equivalent to the Java GraphQL runtime in breadth. It maintains a static registry of supported operations and dispatches repository-backed queries and mutations through the Rust service layer.

The runtime supports:
\begin{itemize}
\item read operations such as unit lookup, search, metadata lookup, health, lock state and relation/association reads,
\item mutation operations such as status changes, unit storage, attribute creation, lock handling and relation/association updates,
\item operation registration and allowlist filtering,
\item GraphQL-style error envelopes for execution failures.
\end{itemize}

This is a compatibility-oriented runtime surface rather than a full clone of the Java GraphQL engine. It exposes a practical operational subset while broader parity is still being developed.

\subsection{Python}

Python serves two different roles:
\begin{description}
\item[Embedded consumer of Rust] Through the \texttt{PyIpto} module exported by Rust via PyO3.
\vspace{.2cm}
\item[Operational tooling] Through scripts that provision test environments, run integration tests, or execute backend benchmarks.
\end{description}

The second role is incidental from an architectural point of view. The first role is structural.

\subsection{Current parity position}

The Rust implementation status:
\begin{itemize}
\item repository core behavior is largely implemented,
\item relations, associations, locking and status-transition surfaces are present,
\item PostgreSQL and Neo4j backends are both active targets,
\item the PyO3 surface is substantial rather than minimal,
\item GraphQL support is present but still partial compared to Java,
\item exact search and lifecycle parity still has edge-case gaps.
\end{itemize}

The Rust implementation is an active parity-seeking runtime, not an experimental side branch. It already carries meaningful repository behavior, a Python-facing integration surface, backend abstraction, and a growing GraphQL compatibility layer.
